\section{Discussion}
\label{sec:discussion}

\paragraph{Limitations.}
The main limitation is computational. The full assignment model (Step~3) and the
three-stage blueprint program (Step~5) embed sequencing constraints with large
Big-$M$ coefficients that weaken the LP relaxation, leaving wide optimality gaps
within the imposed time limits (300\,s and 600\,s). This is not merely a run-time
issue: it produces the spurious negative Step-3 VSS
(Section~\ref{subsec:VSS}) and means the operational cost of the blueprint can
only be bounded to within the residual gap rather than measured exactly
(Section~\ref{sec:bp_apply}). The blueprint is also built from a small scenario
tree and an untuned penalty $\beta_5 = 1$, so the recommended quotas reflect the
sampled waiting lists rather than a calibrated risk preference, and its scope is a
single specialty (ORT). Finally, durations are sampled independently per patient,
ignoring day-level correlation, and the worst-case-tight provisioning leaves no
headroom---as the infeasible comparison seed shows.

\paragraph{Further research.}
A sensitivity sweep over $\beta_5$ would expose the full trade-off between
blueprint tightness and operational flexibility. Tightening the
formulation---logic-based or indicator constraints in place of Big-$M$, valid
inequalities, or a Benders decomposition exploiting the scenario-separable second
stage---would close the gap and turn the blueprint comparison into an exact
measurement. Larger and correlated scenario models, a small safety margin on the
dominant subgroup, and an extension to a multi-specialty, multi-week horizon with
a risk-averse objective (e.g.\ CVaR on overtime) would bring the model closer to
tactical master-schedule practice.
